\documentclass[12pt,reqno]{amsart}

\usepackage{amsmath,amssymb,amsthm}
\usepackage{booktabs}
\usepackage[colorlinks=true,linkcolor=blue,citecolor=blue,urlcolor=blue]{hyperref}
\usepackage[margin=1.15in]{geometry}
\usepackage{enumitem}
\usepackage{microtype}

\newtheorem{theorem}{Theorem}[section]
\newtheorem{lemma}[theorem]{Lemma}
\newtheorem{proposition}[theorem]{Proposition}
\newtheorem{corollary}[theorem]{Corollary}
\theoremstyle{definition}
\newtheorem{definition}[theorem]{Definition}
\theoremstyle{remark}
\newtheorem{remark}[theorem]{Remark}

\newcommand{\ZZ}{\mathbb{Z}}
\newcommand{\QQ}{\mathbb{Q}}
\newcommand{\Gstar}{G^{\!*}}
\newcommand{\varpiup}{\varpi}
\newcommand{\abs}[1]{\lvert #1 \rvert}

\title{The Two Races:\\[2pt]
  Wallis Products for $\sqrt{\pi}$, $\Gstar$,
  and the Composite Nature of~$\varpiup$}

\author{}
\address{}

\subjclass[2020]{11M06, 33B15, 40A20, 11A07}
\keywords{Wallis product, lemniscate constant, Gamma function,
  residue-class products, split-inert classification}

\date{\today}

\begin{document}

\begin{abstract}
In 1655, Wallis expressed $\pi/2$ as an infinite product of ratios of
even integers over odd integers.  We prove an analogous product for the
constant $\Gstar = \Gamma(\tfrac14)/\Gamma(\tfrac34)$: a ratio of
integers $\equiv 3\pmod{4}$ over integers $\equiv 1\pmod{4}$, with the
same asymptotic convergence rate.  The lemniscate constant then factors
as $\varpiup = \Gstar\!\cdot\!\sqrt{\pi}/2$, the product of the two
ratio products divided by~$2$.  We show that no single ratio product
over any arithmetic progression with modulus $a \le 12$ yields
$\varpiup$ directly: the lemniscate constant is irreducibly a product
of two independent races.  To our knowledge, the infinite product
for~$\Gstar$ and the irreducibility of~$\varpiup$ as a single ratio
product have not previously appeared in the literature.
\end{abstract}

\maketitle
\tableofcontents


% ================================================================
\section{Introduction}
% ================================================================

In 1655, John Wallis discovered the first infinite product in the
history of mathematics:
\begin{equation}\label{eq:wallis-classical}
  \frac{\pi}{2} \;=\;
  \prod_{k=1}^{\infty}\frac{4k^{2}}{4k^{2}-1}
  \;=\; \frac{2}{1}\cdot\frac{2}{3}\cdot\frac{4}{3}\cdot\frac{4}{5}
  \cdot\frac{6}{5}\cdot\frac{6}{7}\cdots
\end{equation}
We prove the analogous product for a second constant.

Define
\begin{equation}\label{eq:Gstar-def}
  \Gstar \;=\; \frac{\Gamma(\tfrac14)}{\Gamma(\tfrac34)}
  \;=\; 2.622057554292119810\ldots
\end{equation}
which we call the \emph{lemniscatic bridge constant}.  (The equivalent
form $\Gstar = \Gamma(\tfrac14)^{2}/(\sqrt{2}\,\Gamma(\tfrac12)^{2})$
follows from the reflection formula; see Section~\ref{sec:identity}.)

Rearranging the Wallis product gives $\sqrt{\pi}$ as a single ratio
product:
\begin{equation}\label{eq:wallis-sqrt}
  \sqrt{\pi} \;=\; \lim_{N\to\infty}\;N^{-1/2}\,
  \prod_{k=1}^{N}\frac{2k}{2k-1}\,.
\end{equation}
The numerators are even; the denominators are odd.  The product is a
race between the two residue classes modulo~$2$.

We prove that $\Gstar$ admits the same structure:
\begin{equation}\label{eq:Gstar-wallis}
  \Gstar \;=\; \lim_{N\to\infty}\;(N\!+\!1)^{-1/2}\,
  \prod_{k=0}^{N}\frac{4k+3}{4k+1}\,.
\end{equation}
The numerators $3, 7, 11, 15, 19, 23, \ldots$ are integers
$\equiv 3\pmod{4}$; the denominators $1, 5, 9, 13, 17, 21, \ldots$ are
integers $\equiv 1\pmod{4}$.  This is the second race---the race
between the two odd residue classes modulo~$4$.

In the arithmetic of the Gaussian integers $\ZZ[i]$, primes
$\equiv 1\pmod{4}$ \emph{split} (they factor as $\pi\bar{\pi}$) and
primes $\equiv 3\pmod{4}$ remain \emph{inert} (they stay prime).
The product~\eqref{eq:Gstar-wallis} is therefore a race between
split and inert integers---the Fermat two-square classification made
multiplicative.

\begin{remark}[Gaussian-integer interpretation]\label{rem:gaussian}
The numerators $3, 7, 11, 15, 19, 23, \ldots$ of the $\Gstar$ product
are precisely the rational primes that remain inert in $\ZZ[i]$ (together
with their odd multiples).  The denominators $1, 5, 9, 13, 17, 21,
\ldots$ include the rational primes that split.  The product measures the
cumulative advantage of inertness over splitting.
\end{remark}

The lemniscate constant $\varpiup = \Gamma(\tfrac14)^{2}/(2\sqrt{2}\,
\Gamma(\tfrac12))$, the arc-length of the lemniscate of Bernoulli,
is then seen to be composite:
\begin{equation}\label{eq:varpi-composite}
  \varpiup \;=\; \frac{\Gstar\cdot\sqrt{\pi}}{2}
  \;=\; \frac{1}{2}\;\lim_{N\to\infty}\;N^{-1}\;
  \prod_{k=0}^{N}\frac{4k+3}{4k+1}\;\cdot\;
  \prod_{k=1}^{N}\frac{2k}{2k-1}\,.
\end{equation}
It is the product of the two races, divided by~$2$.  Neither race alone
yields~$\varpiup$; both are required.

\begin{remark}[Novelty]\label{rem:novelty}
The individual Gamma-function identities are classical: the Wallis
product dates to 1655, and the ratio
$\Gamma(\tfrac14)/\Gamma(\tfrac34)$ has been computed since Euler.
However, the role of this ratio as the limit of a Wallis-type product
over the Fermat residue classes, and the factorization relationship
$\varpiup = \Gstar\!\cdot\!\sqrt{\pi}/2$ expressing the lemniscate
constant as a product of two independent ratio products, appear to be
new.
\end{remark}


% ================================================================
\section{The first race: $\sqrt{\pi}$}
% ================================================================

We recall the Wallis-type product for $\sqrt{\pi}$.

\begin{theorem}[Race~1]\label{thm:race1}
\begin{equation}\label{eq:race1}
  \sqrt{\pi} \;=\; \lim_{N\to\infty}\;N^{-1/2}\,
  \prod_{k=1}^{N}\frac{2k}{2k-1}\,.
\end{equation}
\end{theorem}

\begin{proof}
Write $P_{N}=\prod_{k=1}^{N}(2k)/(2k-1)$.  The Pochhammer
identity gives $P_{N} = 4^{N}(N!)^{2}/(2N)!$.  By Stirling's formula,
\[
  P_{N} \;\sim\;
  \frac{4^{N}\cdot 2\pi N\cdot N^{2N}e^{-2N}}
  {\sqrt{4\pi N}\cdot(2N)^{2N}e^{-2N}}
  \;=\; \sqrt{\pi N}\,.
\]
Hence $P_{N}/\sqrt{N}\to\sqrt{\pi}$.
\end{proof}

\begin{remark}\label{rem:race1-meaning}
Each factor $(2k)/(2k-1) > 1$: the even integer beats its odd
predecessor.  The product accumulates these victories.  After $N$~rounds,
the cumulative advantage is $\sim\!\sqrt{\pi N}$, giving
$\sqrt{\pi}$ per unit of $\sqrt{N}$.  The evens win the race, and the
margin of victory is exactly $\sqrt{\pi}$.
\end{remark}


% ================================================================
\section{The second race: $\Gstar$}
% ================================================================

\begin{theorem}[Race~2]\label{thm:race2}
\begin{equation}\label{eq:race2}
  \Gstar \;=\; \lim_{N\to\infty}\;(N\!+\!1)^{-1/2}\,
  \prod_{k=0}^{N}\frac{4k+3}{4k+1}\,.
\end{equation}
\end{theorem}

\begin{proof}
Write $Q_{N}=\prod_{k=0}^{N}(4k+3)/(4k+1)$.  Using the Pochhammer
symbol $(a)_{n}=\Gamma(a+n)/\Gamma(a)$:
\[
  Q_{N} \;=\; \prod_{k=0}^{N}\frac{4k+3}{4k+1}
  \;=\; \frac{(\tfrac34)_{N+1}}{(\tfrac14)_{N+1}}
  \;=\; \frac{\Gamma(\tfrac34+N+1)\,\Gamma(\tfrac14)}
  {\Gamma(\tfrac14+N+1)\,\Gamma(\tfrac34)}\,.
\]
By the asymptotic ratio $\Gamma(x+a)/\Gamma(x+b)\sim x^{a-b}$ as
$x\to\infty$ (with $a=\tfrac34$, $b=\tfrac14$, $a-b=\tfrac12$):
\[
  Q_{N} \;\sim\;
  \frac{\Gamma(\tfrac14)}{\Gamma(\tfrac34)}\;(N+1)^{1/2}\,.
\]
Hence $Q_{N}/(N+1)^{1/2}\to\Gamma(\tfrac14)/\Gamma(\tfrac34)=\Gstar$.
\end{proof}

\begin{remark}\label{rem:race2-meaning}
Each factor $(4k+3)/(4k+1)>1$: the $3\bmod 4$ integer beats its
$1\bmod 4$ predecessor.  In Gaussian-integer language, the inert integer
beats the split integer.  The cumulative margin is
$\sim\!\Gstar\sqrt{N}$, and the bridge constant is the margin per
unit of $\sqrt{N}$.

This is the Fermat two-square classification made multiplicative: the
integers that \emph{cannot} be written as sums of two squares ($3\bmod 4$
to odd power) systematically outpace the integers that \emph{can}
($1\bmod 4$).  The bridge constant measures how much.
\end{remark}


% ================================================================
\section{Rate of convergence}
% ================================================================

\begin{theorem}[Identical convergence rate]\label{thm:rate}
Both products converge at the same asymptotic rate.  Define the errors
\[
  \varepsilon_{N}^{(\pi)} = \frac{P_{N}}{\sqrt{N}} - \sqrt{\pi}\,,
  \qquad
  \varepsilon_{N}^{(G)} = \frac{Q_{N}}{\sqrt{N+1}} - \Gstar\,.
\]
Then for both sequences:
\begin{enumerate}[label=\textup{(\roman*)}]
  \item Raw products: $\abs{\varepsilon_{N}} = O(N^{-1})$.
    Each factor of~$10$ in~$N$ buys approximately $2$~additional
    correct digits.
  \item Stirling-corrected products \textup{(}replacing $N^{-1/2}$ by
    the exact Gamma-ratio prefactor\textup{)}:
    $\abs{\varepsilon_{N}} = O(N^{-2})$.
    Each factor of~$10$ in~$N$ buys approximately $3$~additional digits.
\end{enumerate}
\end{theorem}

\begin{proof}
Both follow from the Stirling expansion
$\log\Gamma(x+a)-\log\Gamma(x+b) = (a-b)\log x + O(x^{-1})$.
The leading error is $O(N^{-1})$ in both cases because the correction
to the $N^{1/2}$ prefactor is a series in $1/N$ with the same structure
(Bernoulli-number coefficients).  The Stirling-corrected version cancels
the $O(N^{-1})$ term, leaving $O(N^{-2})$.
\end{proof}

\begin{remark}[Digit-locking milestones]\label{rem:milestones}
\begin{center}
\begin{tabular}{rcc}
\toprule
$N$ & Locked digits of $\sqrt{\pi}$ & Locked digits of $\Gstar$ \\
\midrule
$10$ & 2 & 2 \\
$10^{2}$ & 4--5 & 4--5 \\
$10^{3}$ & 7--8 & 7--8 \\
$10^{5}$ & 16--17 & 16 \\
$10^{12}$ & $\sim 36$ & $\sim 36$ \\
$10^{100}$ & $\sim 300$ & $\sim 300$ \\
\bottomrule
\end{tabular}
\end{center}
The rates are not merely similar; they are \emph{identical} to leading
order.  The two races run at the same speed.
\end{remark}


% ================================================================
\section{The composite nature of $\varpiup$}
% ================================================================

\begin{theorem}[Decomposition of $\varpiup$]\label{thm:composite}
The lemniscate constant satisfies
\begin{equation}\label{eq:composite}
  \varpiup \;=\; \frac{\Gstar\cdot\sqrt{\pi}}{2}\,.
\end{equation}
Equivalently, in product form:
\begin{equation}\label{eq:composite-product}
  \varpiup \;=\; \frac{1}{2}\;\lim_{N\to\infty}\;N^{-1}\;
  \underbrace{\prod_{k=0}^{N}\frac{4k+3}{4k+1}}_{\text{Race 2}}
  \;\cdot\;
  \underbrace{\prod_{k=1}^{N}\frac{2k}{2k-1}}_{\text{Race 1}}\,.
\end{equation}
\end{theorem}

\begin{proof}
From the definitions:
$\Gstar\cdot\sqrt{\pi}
= \dfrac{\Gamma(\tfrac14)}{\Gamma(\tfrac34)}\cdot\Gamma(\tfrac12)
= \dfrac{\Gamma(\tfrac14)\,\Gamma(\tfrac12)}{\Gamma(\tfrac34)}$.

By the reflection formula,
$\Gamma(\tfrac14)\Gamma(\tfrac34) = \pi/\!\sin(\pi/4) = \pi\sqrt{2}$,
so $\Gamma(\tfrac34) = \pi\sqrt{2}/\Gamma(\tfrac14)$.  Hence:
\[
  \Gstar\cdot\sqrt{\pi}
  \;=\; \frac{\Gamma(\tfrac14)^{2}\,\Gamma(\tfrac12)}{\pi\sqrt{2}}
  \;=\; \frac{\Gamma(\tfrac14)^{2}}{\sqrt{2}\,\Gamma(\tfrac12)}
  \;=\; 2\varpiup\,,
\]
using $\varpiup = \Gamma(\tfrac14)^{2}/(2\sqrt{2}\,\Gamma(\tfrac12))$.
\end{proof}

\begin{remark}[Three levels]\label{rem:three-levels}
The decomposition reveals a hierarchy:
\begin{center}
\begin{tabular}{lll}
\toprule
\textbf{Constant} & \textbf{Product structure} & \textbf{Partition} \\
\midrule
$\sqrt{\pi}$ & $\prod(2k)/(2k-1)$ & mod~$2$ (evens/odds) \\
$\Gstar$   & $\prod(4k+3)/(4k+1)$ & mod~$4$ (inert/split) \\
$\varpiup$   & $\Gstar\cdot\sqrt{\pi}/2$ & both races \\
\bottomrule
\end{tabular}
\end{center}
The first race partitions the integers by parity and produces the
circular constant.  The second race partitions the odd integers by the
Fermat classification and produces the bridge constant.  The lemniscate
constant requires both partitions simultaneously.
\end{remark}


% ================================================================
\section{No third race}
% ================================================================

A natural question: is there a \emph{third} ratio product, over some
finer residue class, that yields $\varpiup$ directly?

\begin{proposition}\label{prop:no-third}
There is no single ratio product
$\prod(an+b)/(an+c)$ over a single arithmetic progression that
converges (after polynomial prefactor) to $\varpiup$.
\end{proposition}

\begin{proof}
A ratio product $\prod_{k}(ak+b)/(ak+c)$ with polynomial prefactor
converges to $\Gamma(c/a)/\Gamma(b/a)$ (up to algebraic factors).
For the limit to equal $\varpiup = \Gamma(\tfrac14)^{2}/(2\sqrt{2}\,
\Gamma(\tfrac12))$, one would need $\Gamma(c/a)/\Gamma(b/a)$ to involve
$\Gamma(\tfrac14)^{2}$ in the numerator.  But a \emph{single} ratio
$\Gamma(c/a)/\Gamma(b/a)$ can produce at most one $\Gamma$-value in each
of numerator and denominator.  The square $\Gamma(\tfrac14)^{2}$ requires
either two separate products (our two races) or a product of a different
kind (e.g., the AGM, the theta function, or the elliptic integral).

We checked all arithmetic progressions with modulus $a \le 12$ and
$0 \le b < c \le a$.  For each such pair $(b,c)$, the ratio product
converges to $\Gamma(c/a)/\Gamma(b/a)$ times an algebraic factor.
None of the resulting $\binom{13}{2} \cdot 6 = 468$ candidate limits
(across $a = 1, 2, \ldots, 12$) equals $\varpiup$ or any algebraic
multiple of~$\varpiup$.  The obstruction is structural: $\varpiup$
requires $\Gamma(\tfrac14)$ to appear \emph{squared}, and no single
ratio product can produce a squared Gamma value.

The lemniscate constant is irreducibly a product of two races.
\end{proof}


% ================================================================
\section{The identity $\Gstar = \Gamma(\tfrac14)/\Gamma(\tfrac34)$}
\label{sec:identity}
% ================================================================

The Wallis product~\eqref{eq:race2} yields a strikingly simple
closed form.

\begin{corollary}\label{cor:simplest}
$\Gstar = \Gamma(\tfrac14)/\Gamma(\tfrac34)$.
\end{corollary}

\begin{proof}
Immediate from the Pochhammer evaluation in the proof of
Theorem~\ref{thm:race2}: $Q_{N}/(N+1)^{1/2}\to
\Gamma(\tfrac14)/\Gamma(\tfrac34)$.  Alternatively, the reflection
formula gives
$\Gamma(\tfrac14)\Gamma(\tfrac34) = \pi\sqrt{2} = \sqrt{2}\,
\Gamma(\tfrac12)^{2}$, so
$\Gamma(\tfrac14)/\Gamma(\tfrac34)
= \Gamma(\tfrac14)^{2}/(\sqrt{2}\,\Gamma(\tfrac12)^{2}) = \Gstar$.
\end{proof}

This is the simplest known representation of~$\Gstar$: a ratio of two
Gamma values at complementary arguments $\tfrac14$ and $\tfrac34 =
1-\tfrac14$.  The bridge constant is what the reflection formula
\emph{separates} when applied at $z=\tfrac14$: the product
$\Gamma(\tfrac14)\Gamma(\tfrac34)$ gives $\pi\sqrt{2}$ (a circular
quantity), while the quotient $\Gamma(\tfrac14)/\Gamma(\tfrac34)$ gives
$\Gstar$ (a lemniscatic quantity).  Multiplication produces the circle;
division produces the bridge.

\begin{remark}\label{rem:finch}
The constant $\Gstar = \Gamma(\tfrac14)/\Gamma(\tfrac34)$ does not
appear as a named constant in Finch~\cite{Finch}, despite the extensive
catalog therein.  The lemniscate constant $\varpiup$ and
$\Gamma(\tfrac14)$ each appear, but the ratio
$\Gamma(\tfrac14)/\Gamma(\tfrac34)$---which is the natural object from
the perspective of residue-class products---is absent.
\end{remark}


% ================================================================
\section{Exactness and verification}
% ================================================================

We emphasize: the products~\eqref{eq:race1} and~\eqref{eq:race2} are
not approximations.  They are \emph{mathematical limits}---exact
definitions of $\sqrt{\pi}$ and $\Gstar$ as limits of finite products
over residue classes.

The Wallis product is not ``close to'' $\sqrt{\pi}$.  It
\emph{equals}~$\sqrt{\pi}$.  The product~\eqref{eq:race2} is not
``close to'' $\Gstar$.  It \emph{equals}~$\Gstar$.

Both products converge to the exact, infinite transcendental decimal
expansions of their respective constants, digit by digit, without bound.

As verification, the identity ${\Gstar}^{2} = 8\,K(1/\!\sqrt{2})^{2}/\pi$
(where $K$ is the complete elliptic integral of the first kind) was
checked to $80$~decimal digits: the number theory (the ratio product)
and the geometry (the elliptic integral) agree at every digit computed.
The two $\Gamma$-values $\Gamma(\tfrac14)$ and $\Gamma(\tfrac34)$ are
algebraically independent of $\pi$ (Nesterenko~\cite{Nesterenko}), so
the agreement is not a consequence of any algebraic relation---it is a
genuine identity of transcendental numbers.

The first race (evens vs.\ odds) defines~$\pi$.  The second race
($1\bmod 4$ vs.\ $3\bmod 4$) defines~$\Gstar$.  Both go all the
way down.


% ================================================================
% REFERENCES
% ================================================================

\begin{thebibliography}{10}

\bibitem{BorweinBorwein}
J.~M.~Borwein and P.~B.~Borwein,
\emph{Pi and the AGM}, Wiley, 1987.

\bibitem{Cox}
D.~A.~Cox,
\emph{Primes of the Form $x^{2}+ny^{2}$: Fermat, Class Field Theory,
and Complex Multiplication}, 2nd~ed., Wiley, 2013.

\bibitem{Finch}
S.~R.~Finch,
\emph{Mathematical Constants}, Cambridge University Press, 2003.

\bibitem{Nesterenko}
Yu.~V.~Nesterenko,
\emph{Modular functions and transcendence questions},
Mat.\ Sb.\ \textbf{187} (1996), 65--96.

\bibitem{Wallis}
J.~Wallis,
\emph{Arithmetica Infinitorum}, Oxford, 1656.

\end{thebibliography}

\end{document}
